\section{Polynomial-size selector encoding}
\label{sec:selector-independent}

The explicit principal-set disjunction has exponentially many branches when
$n$ varies.  Introduce variables $b_i$ satisfying
$b_i^2-b_i=0$ and put $S_b=\operatorname{diag}(b_1,\ldots,b_n)$.  If
$I=\{i:b_i=1\}$, simultaneously permute rows and columns so that $I$ comes
first.  The matrix
\[
 I-S_b-S_bAS_b
\]
then has block form
\[
 \begin{pmatrix}-A[I,I]&0\\0&I_{I^c}\end{pmatrix}.
\]
Consequently
\begin{equation}
 \det(I-S_b-S_bAS_b)=(-1)^{|I|}\det A[I,I].
 \label{eq:selector-independent}
\end{equation}
The empty selector gives determinant $1$, so it cannot satisfy a strict
negative inequality.

A determinant of an $n\times n$ matrix has a division-free arithmetic circuit
of polynomial size, as do characteristic coefficients and the Hurwitz
minors.  Rational row reduction computes a basis $B$ of
$\operatorname{im}\Gamma$ and a left inverse $C$ with polynomial encoding
length: all numerators and denominators are ratios of minors of the rational
input matrix.  Replacing every arithmetic gate by a fresh variable and one
polynomial equation gives a sparse existential formula of polynomial total
size.  Positivity and negativity are legitimate atomic relations in
$\exists\mathbb R$; alternatively, they can be encoded by
\[
 g>0\iff\exists z\ (gz^2-1=0),\qquad
 g<0\iff\exists z\ ((-g)z^2-1=0).
\]
Thus the all-mobile decision problem belongs to the existential theory of the
reals.  This is only a membership result; no $\exists\mathbb R$-hardness claim
is used.
